% !TeX root = ../main.tex

\chapter{Related Work}\label{chapter:related-work}

The challenge of combining high agility with aerodynamic efficiency defines a long-standing trade-off in UAV design. Pure multirotors provide full six-degree-of-freedom control and extreme maneuverability but waste energy by tilting thrust for lift during forward flight. Fixed-wing UAVs achieve order-of-magnitude better cruise efficiency through aerodynamic lift but lose the ability to hover or hold arbitrary attitudes. Hybrid VTOL platforms attempt to bridge these regimes at the cost of additional actuators, transition logic, and increased control complexity.

\section{Control versus Efficiency Trade-off}

Multirotor control has matured through geometric and Incremental Nonlinear Dynamic Inversion (INDI) frameworks that achieve aggressive trajectory tracking without detailed modeling~\cite{Sieberling2010,Goodarzi2015,Tzoumanikas2021}. These architectures rely solely on measured accelerations and remain robust under aerodynamic disturbances. However, they operate under the implicit assumption that thrust is the sole lift source---an assumption that breaks once fixed lifting surfaces contribute significant forces. The question is therefore whether these established control laws remain valid when lift is partly aerodynamic and unmodeled.

\section{Existing Approaches and Their Limits}

Hybrid VTOL concepts such as tilt-rotor and tailsitter aircraft achieve efficient cruise flight and transitions~\cite{Tal2022Global,DaudFilho2024}, yet their complexity is high. They require coordinated mode switching, non-linear actuator coupling, and precise aerodynamic modeling---elements that obscure the fundamental control question: can a standard multirotor controller tolerate significant aerodynamic forces without explicit modeling?

\section{Aerodynamic Augmentation Without Added Complexity}

Prior studies have explored fixed aerodynamic surfaces on multirotors as a simpler alternative. Dawkins and DeVries~\cite{Dawkins2018} demonstrated $\approx 35\,\%$ lower cruise power at 3--5\,\si{\meter\per\second} with small attached wings, while Xiao et al.~\cite{Xiao2020} reported roughly $50\,\%$ electrical-power reduction at 15\,\si{\meter\per\second} on a 1.2\,kg platform, both without altering the baseline controller. These results suggest that passive aerodynamic surfaces can meaningfully reduce power consumption while maintaining full control authority. However, none of these works analyzed how unmodeled aerodynamic forces affect control stability or trajectory-tracking accuracy---especially under a disturbance-rejection scheme such as INDI.

\section{Research Rationale}

The literature therefore leaves an open gap: no prior study has systematically tested whether a disturbance-rejection controller such as INDI can maintain precision flight on a multirotor with strong, yet unmodeled, aerodynamic lift. Addressing this gap requires a configuration that introduces aerodynamic coupling without adding mechanical or algorithmic complexity. A quadrotor with fixed wings---an aerodynamic surface-enhanced multirotor---satisfies that criterion.

Consequently, this thesis focuses on (i) implementing an INDI-based controller on such a platform, (ii) quantifying its tracking accuracy under aerodynamic loading, and (iii) evaluating whether explicit aerodynamic modeling is unnecessary within the tested speed envelope.

\section{Summary and Transition}

Existing research establishes the energy-efficiency benefits of fixed aerodynamic surfaces on multicopters but leaves their influence on control robustness largely unexplored. The missing link is a quantitative examination of how a conventional multirotor controller---specifically, an INDI-based architecture---behaves when aerodynamic lift and drag become significant yet remain unmodeled.

To address this, the present work isolates the essential physics by selecting the simplest possible testbed: a standard quadrotor geometry augmented with symmetric fixed wings. The next chapter formalizes its six-degree-of-freedom dynamics and defines the aerodynamic forces that the controller will deliberately ignore. This provides the analytical baseline against which the later experimental results can be interpreted.
